\documentclass[11pt]{article}

\usepackage[margin=1in]{geometry}
\usepackage{amsmath,amssymb}
\usepackage{booktabs}
\usepackage{siunitx}
\usepackage{xcolor}
\usepackage{hyperref}
\hypersetup{colorlinks=true,linkcolor=blue,urlcolor=blue}

\newcommand{\Astored}{A_{\mathrm{stored}}}
\newcommand{\Imeas}{I_{\mathrm{measured}}}
\newcommand{\Pphys}{P_{\mathrm{physical}}}
\newcommand{\Ptrain}{P_{\mathrm{training}}}
\newcommand{\Psim}{P_{\mathrm{sim}}}
\newcommand{\code}[1]{\texttt{#1}}

\title{How PtychoPINN--CI Defines Its Diffraction Target,\\
and Why the Stale--Probe Bug Was Invisible to It}
\author{CI compatibility ablation --- session note}
\date{2026-07-15}

\begin{document}
\maketitle

\begin{abstract}
The count--intensity (CI) contract never imports or measures a physically
calibrated diffraction target. It \emph{constructs} the target from the stored
normalized amplitudes together with a single scalar it derives from those same
amplitudes. This makes CI's notion of ``physical scale'' entirely
self--referential. We show that the derived scalar cancels out of the
object--recovery problem, so the only quantity in the entire chain that
carries the identity of the physical illumination is the stored probe. This is
the structural reason the stale--probe bug --- serializing the pre-normalization
probe $Q$ while the simulator generated data with $kQ$ --- was undetectable
from the target alone, and why it had to be fixed at the probe.
All claims are grounded in the code on branch
\code{feature/ci-compatibility-ablation}.
\end{abstract}

\section{Setup and question}

Let $\Astored$ denote the stored, normalized--amplitude diffraction written to
the dataset NPZ, and let $R$ denote the stored probe (either the legacy
\code{probeGuess} or the simulation probe \code{probe\_simulated}). The question
this note answers:

\begin{quote}
\emph{How does CI know what the ``target'' (physically scaled) diffraction
values are?}
\end{quote}

The short answer: it does not know them independently. It builds them.

\section{CI constructs the target from stored data}

The CI adapter converts normalized amplitudes to count--domain intensities in
\code{ptycho\_torch/scaling\_contract.py}. The operative lines are
\begin{verbatim}
scale          = count_amplitude_scale         # = S
intensity      = (scale * amplitude).square()  # I_measured = (S * A_stored)^2
probe_physical = scale * probe                 # P_physical  = S * R
\end{verbatim}
so that, with $S \equiv \code{count\_amplitude\_scale}$,
\begin{align}
\Imeas &= \bigl(S\,\Astored\bigr)^2, \label{eq:target}\\
\Pphys &= S\,R. \label{eq:pphys}
\end{align}

The scalar $S$ is itself \emph{derived from the stored amplitudes}, via
\code{derive\_intensity\_scale\_from\_amplitudes} (called at
\code{scripts/studies/grid\_lines\_torch\_runner.py:1524}). It is chosen to map
the diffraction energy onto a fixed geometric target. The native dataloader
performs the equivalent normalization directly
(\code{ptycho\_torch/dataloader.py:1264--1266}):
\begin{equation}
\code{target\_energy} = \Bigl(\tfrac{N}{2}\Bigr)^2,
\qquad
\code{rms\_input\_scale} = \sqrt{\frac{\code{target\_energy}}{\langle
\text{measured}^2\rangle}} ,
\label{eq:rms}
\end{equation}
with, for $N=128$, $\code{target\_energy} = 64^2 = 4096$. The CI Poisson target
tensor \code{measured\_intensity} is a verbatim copy of the stored diffraction
stack (\code{dataloader.py:1160--1162}); no external photon budget is consulted.

Empirically, for the lines/$N128$ bundle used throughout the ablation,
$S = \code{count\_amplitude\_scale\_test} \approx 329.2$
(recorded in the postfix re-evaluation metrics).

\paragraph{Consequence.} $S$ is a pure numerical--conditioning gauge computed
from the data to hit a canonical RMS. It carries no physical calibration of its
own; the absolute scale of $\Imeas$ is inherited entirely from $\Astored$,
i.e.\ from whatever was written to the NPZ at simulation time.

\section{The derived scale $S$ cancels}

The physics forward model predicts intensity by propagating the network texture
$z$ through the physical probe:
\begin{equation}
\lambda(z) \;=\; \bigl|\mathcal{F}\{\Pphys\, z\}\bigr|^2
\;=\; \bigl|\mathcal{F}\{S R z\}\bigr|^2
\;=\; S^2\,\bigl|\mathcal{F}\{R z\}\bigr|^2 .
\end{equation}
Matching the prediction to the target~\eqref{eq:target},
\begin{equation}
S^2\,\bigl|\mathcal{F}\{R z\}\bigr|^2
\;=\; \bigl(S\,\Astored\bigr)^2
\;=\; S^2\,\Astored^2
\quad\Longrightarrow\quad
\boxed{\;\bigl|\mathcal{F}\{R z\}\bigr|^2 = \Astored^2\;}.
\label{eq:cancel}
\end{equation}
$S$ vanishes identically. The object $z$ that satisfies the fit depends only on
$R$ (the probe) and $\Astored$ (the data), never on $S$.

This is why inspecting the network--facing probe was misleading. The loader
supplies a normalized $\Ptrain = \Pphys / n$ and an \code{output\_scale} factor
with $\code{output\_scale}\cdot\Ptrain = \Pphys$
(audit identity CI--SCALE--001, \code{ptycho\_torch/model.py:1591}). The
normalized $\Ptrain$ is scale--invariant --- identical for $Q$ and $kQ$ --- but
the reconstituted effective probe $\Pphys = S R$ is not, because
\code{output\_scale} restores the $R$--dependent magnitude.

\section{$R$ is the sole carrier of physical identity}

Equation~\eqref{eq:cancel} is internally self--consistent for \emph{any} probe
$R$: the network is always free to move $z$ so that
$|\mathcal{F}\{R z\}|^2 = \Astored^2$. Nothing in the target constrains $R$,
because the target was built from $\Astored$ and a self--derived $S$ that
cancels. Therefore the identity of the physical illumination --- the probe that
actually formed the recorded counts --- enters the problem \emph{only} through
$R$, the array the dataset writer chose to serialize.

\section{The stale--probe bug}

\subsection{Simulation silently rescaled the probe}

Let $Q$ be the extrapolated complex \code{probeGuess} ($128\times128$).
\code{set\_probe\_guess}$(Q)$ forwards to \code{set\_probe}
(\code{ptycho/probe.py:52--66}), which computes a mask--weighted mean amplitude
with the radius--$N/4$ mask $M$,
\begin{equation}
d = \code{probe\_scale}\cdot\operatorname{mean}\bigl|M\odot Q\bigr|
  = 0.3156714,
\end{equation}
and installs the \emph{divided} probe in global simulation state:
\begin{equation}
\Psim = \frac{Q}{d} = k\,Q, \qquad k = 3.167851 .
\end{equation}
The mask only sets the scalar $d$; the full complex field is divided by it.
Hence $Q$ and $\Psim$ share spatial structure and phase and differ only in
amplitude.

\subsection{Simulator and writer used different variables}

\code{generate\_data()} read the global \code{params["probe"]}, so the
measurements were formed with $\Psim = 3.167851\,Q$. The old writer then
serialized its \emph{unchanged local} $Q$ as \code{probeGuess}. The NPZ
therefore implicitly claimed the counts were made with $Q$, when they were made
with $kQ$ --- an amplitude error of $k$ and an intensity error of
\begin{equation}
k^2 = 10.0353 .
\end{equation}

\subsection{CI faithfully propagated the mistake}

Substituting the two probes into~\eqref{eq:cancel} with the data
$\Astored^2 = |\mathcal{F}\{kQ\,z_{\mathrm{true}}\}|^2$ fixed:
\begin{center}
\begin{tabular}{lll}
\toprule
& \textbf{Stale } $R=Q$ & \textbf{Correct } $R=kQ$ \\
\midrule
Normalized $\Ptrain$ mean & $0.25$ & $0.25$ \\
\code{output\_scale}      & $132.24$ & $418.90$ \\
Effective--probe mean     & $33.06$  & $104.73$ \\
Recovered texture         & $z = k\,z_{\mathrm{true}}$ & $z = z_{\mathrm{true}}$ \\
\bottomrule
\end{tabular}
\end{center}
With the stale probe the fit still closes --- but only by inflating the object
by $k$. This is why the failure first looked like a benign object--probe gauge
ambiguity. Empirically the stale H1 texture amplitude was $9.328$ versus
corrected $2.979$ (ground truth $\approx 2.97$); the ratio $3.13 \approx k$.

\section{Why it was invisible, and why it still hurt}

\paragraph{Invisible.} Every check that inspects the \emph{target} or the
\emph{normalized} probe is blind to the bug: the target~\eqref{eq:target} is
self--consistent for any $R$, and $\Ptrain$ is $R$--scale--invariant. The
discriminating quantity is $\Pphys = SR$, whose value is fixed solely by which
probe was serialized. In \code{compute\_loss} the un--normalized
\code{probe\_physical} is merely shape/dtype--audited
(\code{model.py:2676--2682}); the effective probe enters as the product
$\code{output\_scale}\cdot\Ptrain$, and \code{output\_scale} is derived from the
stored probe's magnitude.

\paragraph{Not a harmless gauge shift.} A global object rescale $z \mapsto kz$
would be invisible only if training were gauge--invariant. It is not. Starting
from identical weights and $s_1,s_2\approx 1$, the Poisson gradient
\begin{equation}
\frac{\partial \ell}{\partial \lambda}
= \frac{1 - y/\lambda}{\bar y}
\end{equation}
changes both $\lambda$ and its Jacobian under $R\mapsto kR$; this is not a
common multiplier that Adam's normalization would absorb. Adam moments,
gradient clipping, normalization layers, and the fixed 20--epoch schedule are
\emph{not} transformed into the equivalent gauge. The stale run entered the
high--amplitude decoder basin and strongly excited the periodic (lattice) mode:
\begin{center}
\begin{tabular}{lcc}
\toprule
& amplitude SSIM & lattice fraction \\
\midrule
Stale H1        & $0.6267$ & $37.0\%$ \\
Corrected Phase E & $0.9335$ & $3.54\%$ \\
\bottomrule
\end{tabular}
\end{center}

\section{The Phase E correction and current-state warning}

Phase E preserves both meanings explicitly:
\begin{itemize}
\item \code{probeGuess} $= Q$: the pre--\code{set\_probe} probe, retained for
legacy compatibility;
\item \code{probe\_simulated} $= kQ$: the actual illumination that generated the
synthetic counts.
\end{itemize}
CI selects \code{probe\_simulated}; legacy paths retain \code{probeGuess}
(producer: E1, \code{9053e9f1f}; consumer/selector: E2, \code{114d797ae}).

\paragraph{Warning.} The shared checkout is currently on \code{df997f3ea}
(branch \code{integration/main-ci-session-port-20260715}, rooted at
\code{main}). That HEAD contains no \code{probe\_simulated} producer, consumer,
or provenance selection. The Phase E correction is therefore \emph{absent} from
the checked--out integration branch, and the stale--probe behavior can recur
there until E1/E2 are ported.

\appendix
\section{Code map}
\begin{center}
\begin{tabular}{ll}
\toprule
Quantity & Location (branch \code{feature/ci-compatibility-ablation}) \\
\midrule
$\Imeas = (S\Astored)^2$, $\Pphys = SR$ & \code{ptycho\_torch/scaling\_contract.py} (conversion) \\
$S$ derivation & \code{grid\_lines\_torch\_runner.py:1524} (\code{derive\_intensity\_scale\_from\_amplitudes}) \\
$(N/2)^2$ energy target, \code{rms\_input\_scale} & \code{ptycho\_torch/dataloader.py:1264--1266} \\
\code{measured\_intensity} $=$ stored diffraction & \code{ptycho\_torch/dataloader.py:1160--1162} \\
$\Ptrain,\ \Pphys,\ \code{probe\_normalization}$ split & \code{ptycho\_torch/dataloader.py:1596--1610} \\
CI branch of \code{compute\_loss} & \code{ptycho\_torch/model.py:2638--2730} \\
\code{output\_scale}$\cdot\Ptrain = \Pphys$ audit & \code{ptycho\_torch/model.py:1591} \\
\code{set\_probe} amplitude normalization $d$ & \code{ptycho/probe.py:52--66} \\
\code{probe\_simulated} producer / consumer & \code{9053e9f1f} (E1) / \code{114d797ae} (E2) \\
\bottomrule
\end{tabular}
\end{center}

\end{document}
